\cleardoublepage

\section{总结与展望}

本章对全文工作进行总结，并结合实验结果分析本文方法仍然存在的不足，最后对后续研究方向进行展望。本文围绕视频多模态大模型推理过程中视觉 token 数量过多、prefill 阶段计算开销较大的问题展开研究，以 VideoLLaMA2 为基础视频大模型框架，将面向图像多模态模型的 TokenPacker 视觉 token 压缩思想迁移到视频场景中，构建了逐帧空间压缩、token 校准与轻量微调相结合的视频高效推理方案\cite{videollama22024,tokenpacker2024}。

\subsection{工作总结}

本文首先分析了视频多模态大模型中的视觉 token 冗余问题。对于基于视觉 Transformer 的视频理解模型而言，视频输入通常需要被拆分为多帧图像，再由视觉编码器将每帧图像表示为大量 patch token。以本文使用的 VideoLLaMA2-7B-16F 为例，在输入 16 帧、单帧空间网格为 $24\times24$ 的设置下，单个视频样本在进入后续时空建模模块前会产生 9216 个帧级空间视觉 token。大量视觉 token 会显著增加语言模型上下文长度，并使 prefill 阶段产生较高的计算和显存开销。因此，如何在尽量保持视频语义信息的同时减少视觉 token 数量，是提升视频多模态大模型推理效率的重要问题。

针对上述问题，本文完成了 TokenPacker 在 VideoLLaMA2 中的结构迁移与接口适配。由于 TokenPacker 原始设计主要面向静态图像多模态模型，其输出维度、压缩位置和后端连接方式不能直接套用于 VideoLLaMA2 的时空建模流程。本文没有替换 VideoLLaMA2 原有的 STC Connector，也没有修改语言模型主体，而是在视觉塔输出与 STC Connector 输入之间引入帧级空间 token 压缩模块。该设计将视频 token 压缩问题拆解为逐帧空间压缩问题，在保留原有时间维结构的同时降低每帧内部的空间 token 数量，从而避免破坏 VideoLLaMA2 已有的时空融合机制。

在具体实现上，本文对 TokenPacker 的输入输出结构进行了调整，使其能够接收来自 CLIP 视觉塔的 single features 和 multi-level features，并输出与 VideoLLaMA2 后续模块兼容的视觉特征维度。本文进一步定义了 scale factor 来控制压缩倍率。当 scale factor 分别为 2、3、4 时，每帧视觉 token 数量由 576 分别压缩至 144、64、36，对应单个视频样本的视觉 token 总数由 9216 分别下降到 2304、1024、576。该流程实现了可控的视觉 token 压缩，并形成了完整的“视频输入、帧采样、视觉特征提取、TokenPacker 压缩、STC Connector 时空融合、LLM 解码”的推理链路。

针对 TokenPacker 直接迁移后出现的准确率下降问题，本文进一步设计了 token 校准与轻量监督微调策略。实验表明，直接将图像场景下的 TokenPacker 权重迁移到 VideoLLaMA2 中会带来明显性能下降，其主要原因在于压缩 token 的分布与原始 VideoLLaMA2 后端模块所适配的视觉 token 分布不一致。为缓解这一问题，本文首先使用无标签视频样本进行 token 校准，使 TokenPacker 输出的压缩 token 与目标压缩表示在特征空间中更加接近；随后在校准结果基础上进行轻量监督微调，使压缩模块进一步适配下游视频问答任务。该流程从“特征分布对齐”和“任务目标适配”两个层面改善了压缩模型的性能。

实验部分从压缩效率、模型准确率、压缩倍率鲁棒性和跨数据集泛化能力等角度验证了本文方法的有效性。在 MVBench 主实验中，未压缩 Baseline 的平均推理延迟为 0.4762 秒；引入 TokenPacker 后，当 scale factor 为 2、3、4 时，平均推理延迟分别下降到 0.3699 秒、0.3431 秒和 0.3344 秒，吞吐率分别提升到 2.7037、2.9143 和 2.9900 samples/s\cite{mvbench2024}。进一步分析表明，视觉 token 压缩主要降低的是 prefill 阶段耗时，说明减少视觉上下文长度能够直接缓解视频多模态大模型的输入编码开销。

在准确率恢复方面，本文实验验证了校准和微调的必要性。以 scale factor 为 2 为例，TokenPacker 直接迁移后的 MVBench 准确率为 33.86\%，经过 token 校准后提升到 41.32\%，进一步轻量微调后提升到 47.01\%。在 scale factor 为 3 和 4 的设置下，校准与微调同样带来了稳定增益。这说明 TokenPacker 迁移到视频大模型后的性能下降并非完全来自压缩结构本身，而是与跨模型、跨任务的分布适配问题密切相关；通过适当的校准和微调，压缩模块能够获得更好的任务表现。

本文还比较了 TokenPacker 与 Mean Pooling 在不同压缩倍率下的表现。实验结果表明，Mean Pooling 在低压缩倍率下具有较强竞争力，scale factor 为 2 时准确率达到 50.47\%，高于同倍率下的 TokenPacker 系列方法。这说明简单均值池化由于更接近原始 token 分布，在低压缩场景下更容易与已有 VideoLLaMA2 后端模块保持适配。然而，随着压缩倍率提高，Mean Pooling 的准确率下降更明显；当 scale factor 从 2 增加到 4 时，其准确率从 50.47\% 下降到 39.37\%。相比之下，经过微调的 TokenPacker 从 47.01\% 下降到 41.88\%，下降幅度更小，并在 scale factor 为 4 时超过 Mean Pooling。这表明可学习视觉 token 压缩方法在高压缩倍率下具有更好的信息聚合能力和鲁棒性潜力。

最后，本文使用 Video-MME 对模型进行跨数据集泛化测试\cite{videomme2025}。实验结果显示，TokenPacker 直接迁移版本在 Video-MME 上的总体准确率为 34.70\%，经过校准与轻量微调后提升到 42.95\%。虽然与未压缩 Baseline 的 52.06\% 仍存在差距，但该结果说明本文提出的适配流程并非只对 MVBench 有效，在未参与校准和微调的新数据集上仍能带来一定性能恢复。综合来看，本文验证了将 TokenPacker 迁移至视频大模型进行视觉 token 压缩的可行性，也揭示了可学习压缩方法在高压缩场景下的潜在优势。

\subsection{存在的不足}

尽管本文方法在推理效率和高压缩鲁棒性方面取得了一定效果，但仍然存在以下不足。

第一，压缩模型的准确率仍未完全恢复到未压缩 Baseline 水平。实验结果表明，即使经过 token 校准和轻量微调，TokenPacker 压缩模型在 MVBench 和 Video-MME 上仍与原始 VideoLLaMA2 存在一定差距。这说明当前逐帧空间压缩仍会造成不可忽略的信息损失，尤其是在需要细粒度空间关系、物体交互和复杂时序推理的任务中，压缩后的视觉 token 可能不足以完整保留原始视觉特征中的关键线索。

第二，本文主要关注帧内空间维度的 token 压缩，尚未充分利用视频在时间维度上的冗余信息。实际视频中相邻帧之间通常存在大量重复内容，而不同时间片段对任务答案的重要性也并不相同。本文为了保持 VideoLLaMA2 原有 STC Connector 的输入结构，选择对每一帧进行独立空间压缩，没有进一步设计时间维动态压缩、关键帧选择或跨帧 token 合并机制。因此，在长视频或强时序推理任务中，本文方法对时间冗余的利用仍然有限。

第三，当前压缩策略采用固定 scale factor，缺少面向不同样本和不同任务的自适应机制。在本文实验中，所有视频样本均使用相同压缩倍率，例如每帧统一保留 144、64 或 36 个 token。然而，不同视频的内容复杂度、运动幅度和问题类型差异较大。对于视觉信息简单或背景冗余较多的视频，可以采用更高压缩倍率；而对于包含细粒度动作、文字、多个目标交互的视频，则可能需要保留更多 token。固定压缩倍率虽然便于实验控制，但难以在所有样本上同时达到最优效率和准确率平衡。

第四，本文使用的校准和微调数据规模仍然有限。由于实验资源和时间限制，本文主要从 MVBench 中构造有限规模的校准子集，并进行轻量级训练。该设置能够验证方法可行性，但不足以充分发挥 TokenPacker 作为可学习压缩模块的能力。从子任务结果可以看到，不同任务上的性能恢复并不均衡，部分任务甚至在微调后仍有波动。这说明后续仍需要更大规模、更高质量的视频指令数据和更系统的训练策略，以进一步提升压缩模块的稳定性。

第五，本文对部署层面的优化仍不充分。当前实验主要从模型结构和 token 数量角度分析推理效率，重点统计了延迟、吞吐率和 prefill 时间等指标。但在实际部署中，视频解码、帧采样、视觉塔前向计算、显存调度、KV cache 管理和批处理策略都会影响端到端推理效率。本文尚未针对这些系统因素进行联合优化，因此距离完整的视频大模型部署方案仍有一定距离。

\subsection{未来工作展望}

未来可以从以下几个方向继续改进本文工作。

首先，可以进一步研究时空联合的视觉 token 压缩方法。本文采用逐帧空间压缩以保持 VideoLLaMA2 原有时空结构，但这一策略没有直接处理跨帧冗余。后续可以在空间压缩基础上引入时间维 token 合并、关键帧选择或片段级摘要 token，使模型同时减少空间冗余和时间冗余。特别是在长视频理解场景中，若能够根据视频内容动态保留关键时间片段，将有助于进一步降低上下文长度并提升长视频推理效率。

其次，可以设计动态 token routing 和自适应 token 保留策略。未来模型可以根据视频内容、问题文本和中间特征自动判断不同区域的重要性，为不同帧、不同空间区域分配不同 token 预算。例如，对于包含关键动作或文字信息的区域保留更多 token，而对静态背景或重复帧进行更强压缩。相比固定 scale factor，自适应压缩能够更精细地平衡效率和性能，也更符合视频理解任务中信息分布不均衡的特点。

再次，可以探索更强的校准目标和训练方式。本文主要通过 token 级特征对齐和轻量监督微调缓解分布偏移问题。后续可以进一步引入多层特征蒸馏、注意力分布对齐、答案一致性约束或教师模型指导，使压缩模块不仅对齐局部 token 表示，也能保留原始模型在推理过程中的语义判断能力。此外，还可以结合更大规模的视频指令数据进行阶段式训练，提高压缩模块在不同任务和数据集上的泛化能力。

此外，可以针对不同视频任务构建任务感知的压缩机制。实验结果表明，动作识别、空间关系推理、目标交互、OCR 和计数等任务对视觉信息的需求并不相同。未来可以根据问题类型或任务类别选择不同的压缩策略，例如对 OCR 和细粒度空间推理任务保留更多局部 token，对全局动作识别和场景概括任务使用更强压缩。任务感知压缩有望减少统一压缩策略带来的性能波动。

最后，可以将视觉 token 压缩与系统级推理优化结合起来。视觉 token 压缩主要降低模型输入上下文长度，而实际视频大模型部署还涉及视觉编码器加速、KV cache 管理、批量推理调度和低比特量化等问题。未来可以将本文方法与高效注意力、模型量化、缓存复用和流式视频处理等技术结合，构建面向真实应用场景的端到端高效视频多模态推理系统。

总体而言，本文围绕视频大模型高效推理中的视觉 token 压缩问题开展了初步研究，完成了 TokenPacker 到 VideoLLaMA2 的迁移适配，并通过实验验证了逐帧空间压缩、token 校准和轻量微调的有效性。虽然当前方法在准确率恢复、自适应压缩和长视频处理方面仍有改进空间，但相关实验表明，可学习视觉 token 压缩在高压缩倍率场景下具有值得进一步探索的潜力。随着视频多模态大模型在真实应用中的输入长度和部署需求不断提升，面向视频场景的高效视觉 token 表示与压缩方法仍将是重要的研究方向。
